% Ordinary mathematical proof of the single source association change.
% Requires amsmath/amssymb and the surrounding document's theorem environments.
\section{One shared reassociation removes a shared source pair}
\label{sec:rotation-source-proof}

The useful source transformation keeps every instruction occurrence and
every register operation.  It changes how one repeated three-part sequence
is parenthesized.  Its relevance to state count is that one of the newly
formed pairs is already represented by an existing decoder state.  The
source result is proved below; the simplest 245-state binary realization
is rejected by an instruction-frame counterexample.

\subsection{The repeated fragment}

In the literal source tree, a representative fragment is
\[
 \begin{array}{rcl}
 363&=&\operatorname{seq}(364,382),\\
 382&=&\operatorname{seq}(383,389).
 \end{array}
\]
Let $A$ denote subtree 364, $B$ subtree 383, and $C$ subtree 389.  The
fragment is $A;(B;C)$.  The existing subtree 75 is
\[
 75=\operatorname{seq}(76,94),
\]
where subtree 76 is structurally identical to $A$ and subtree 94 is
structurally identical to $B$.  Therefore $A;B$ already has a decoder.
Replace the shared fragment by
\[
 \boxed{\qquad (A;B);C.\qquad}
\]

The old fragment has nine occurrences.  The following are original
\texttt{Prog254.v} node numbers and their complete paths after the
published constant prefix has been removed.
\begin{center}
\small
\begin{tabular}{rll}
Fragment node & Root-to-fragment path & Old right-child node\\
\hline
363  & \texttt{111010010010100110} & 382\\
1400 & \texttt{11101001001010101111} & 1419\\
1995 & \texttt{11101001001010111011} & 2014\\
2710 & \texttt{11101001001011001001} & 2729\\
2977 & \texttt{111010010010110011110} & 2996\\
4313 & \texttt{111010010010111010100} & 4332\\
4364 & \texttt{111010010010111010101} & 4383\\
4416 & \texttt{111010010010111010110} & 4435\\
4879 & \texttt{1110100100101111001010} & 4898
\end{tabular}
\end{center}

All nine fragment roots represent one shared subtree, and all nine old
right children represent one shared subtree.  Every occurrence is changed.
The latter shared subtree then disappears, while the former becomes a
pair of two existing subtrees.  Exact structural hashing therefore gives
\[
 203\quad\longrightarrow\quad202
\]
distinct internal pair subtrees.  The literal unrolled tree still has
6,459 internal nodes and 6,460 instruction leaves.  Its maximum depth is
still 34 and its maximum consecutive zero run is still seven.

\subsection{Why the reassociation preserves control}

Ordinary associativity of sequential execution is not by itself enough
for this representation: source instructions contain relative break and
continue counts, and zero decrement skips a control subtree.  These counts
are checked against the exact changed paths.

Let $p$ be one of the nine fragment paths and let $s$ be a suffix inside
the indicated component.  The leaf-path correspondence is
\[
 \begin{array}{rcll}
 p0s  &\longmapsto&p00s,&\text{inside }A,\\
 p10s &\longmapsto&p01s,&\text{inside }B,\\
 p11s &\longmapsto&p1s,&\text{inside }C.
 \end{array}
\]
All other leaf paths are unchanged.  This is a bijection preserving leaf
order and the literal instruction attached to each leaf.  In particular,
every register index, every break/continue count, and the terminal leaf
remain the same.

For an instruction at leaf $i$ and a guard $g$ (zero or positive for a
decrement), write $\sigma(i,g)$ for the successor leaf obtained from the
exact program-counter operations and the zero-padded decoder.  The
independent reconstruction verifies the finite identity
\[
 \Phi\bigl(\sigma_P(i,g)\bigr)
 =\sigma_{P'}\bigl(\Phi(i),g\bigr)
\]
at all 8,391 guarded cases, including both decrement branches and all
break/continue cases.  It interprets carry by retaining all cleared carry
bits; it checks that no unconsumed one is discarded at a decoded leaf.
The initial empty PC selects corresponding first leaves.  No unreachable
guard is exempted from this comparison.

\paragraph{Theorem.}
For every natural-valued register vector $R$, the original program $P$
and the rotated program $P'$ have identical instruction traces under
$\Phi$ from their empty-PC roots.  Consequently
\[
 \operatorname{Halt}(P,[],R)
 \quad\longleftrightarrow\quad
 \operatorname{Halt}(P',[],R).
\]
The same register vector appears at corresponding points, and the number
of register instructions before any corresponding point is equal.

\paragraph{Proof.}
At the initial leaves the register files agree.  At corresponding leaves,
the attached operation and its register argument are identical.  Therefore
the two executions select the same decrement guard and apply the same
register update.  The checked successor identity places them at
corresponding next leaves.  Induction on the instruction count proves the
trace claim.  The unique halt is retained, so a finite halting trace in
either program gives one in the other with the same number of register
instructions.  No soundness invariant on formula data, consistency
assumption, or bound on register values is used.\hfill$\square$

\subsection{A rejected 245-state binary realization}

In the original numerical TM labels, the four relevant passive decoder
states represent the following source subtrees:
\[
 363\leftrightarrow178,\qquad
 382\leftrightarrow208,\qquad
 75\leftrightarrow16,\qquad
 389\leftrightarrow206.
\]
The old rows are
\[
\begin{array}{c|cc}
 &0&1\\\hline
178&(0,R,23)&(1,R,208)\\
208&(0,R,24)&(1,R,206)\\
16 &(0,R,23)&(1,R,24).
\end{array}
\]
Change only the two rows of state 178 to
\[
 \delta'(178,0)=(0,R,16),\qquad
 \delta'(178,1)=(1,R,206),
\]
and remove state 208.  Its only old incoming transition is the replaced
read-one transition of state 178.  Every displayed decoder action reads
and writes the same bit and moves right.  The new state 178 thus decodes
the existing $A;B$ subtree on zero and $C$ on one.  All other surviving
rows are unchanged.

This edit produces a literal table with 245 working states, but that
table does not satisfy the inherited instruction-frame contract.  State
178 is also borrowed by a backward decoder alias:
\[
 \delta(218,0)=(0,L,154),\qquad
 \delta(154,0)=(0,R,178).
\]
On the relevant previous-bit-zero context, this path expects the old
read-zero derivative $\delta(178,0)=(0,R,23)$.  The reassociation changes
that derivative to $(0,R,16)$, corrupting other source occurrences.
An explicit failing leaf is
\begin{quote}\ttfamily\small
11101001001010100111000000,
\end{quote}
whose instruction is logical-register-1 decrement.  The candidate can
instead select logical register 2.  The independently checked root phase
invariant does not make these registers equal or force both to zero at
this site.  Thus the finite failure cannot be discarded as an impossible
source branch using that invariant.

The source theorem remains valid.  It does not prove that this binary
table preserves the complete source execution, and the table is not a
certified sub-246 solution.  A successful compiler must charge and check
both ordinary source decoding and all backward alias uses.  Separate
full-frame instruction certificates and a safe blank initialization are
necessary before the end-to-end consistency theorem may be asserted.

\subsection{Independent reconstruction}

\texttt{rotation245\_source\_audit.py} reconstructs the change directly
from the original source node numbers, without reading the searcher's
hash-consed graph.  It parses every source token, checks the complete
forward tree, generates the nine rotations, and uses a direct tree walk
and Boolean-list control operations to compare all guarded successors.
This differs from the searcher's interval/binary-search decoder.  The
certificate includes every old/new leaf path and every guarded successor.
The exact input source SHA-256 is
\begin{quote}\ttfamily\small
1fbc30fcf82a1442ed88d1a1dde412630a431d32a7265d22bdf8517bffe01746.
\end{quote}
The proof above gives the ordinary meaning of those finite checks.  No
new Lean or Rocq run is claimed for the changed source tree by this audit.
